% ============================================================
%  CHƯƠNG I: TỔNG QUAN
% ============================================================
\chapter{TỔNG QUAN}
\label{chap:tong_quan}

% ────────────────────────────────────────────────────────────
\section{Đặt vấn đề}
\label{sec:dat_van_de}

Trong những năm gần đây, hành vi tiêu thụ nội dung giải trí của người dùng tại Việt Nam
và trên toàn cầu đã dịch chuyển mạnh mẽ từ truyền hình truyền thống sang các nền tảng phát
trực tuyến (streaming) trên thiết bị di động. Sự phổ biến của các dịch vụ như Netflix,
Disney+, FPT Play hay Galaxy Play cho thấy mô hình xem phim theo nhu cầu (on-demand) đã
trở thành phương thức tiêu thụ nội dung giải trí chủ đạo, không còn là xu hướng nhất
thời, với thời lượng xem trực tuyến tăng liên tục từng năm và phần lớn lượt xem đến từ
thiết bị di động.

Đằng sau giao diện xem phim đơn giản là những bài toán kỹ thuật phức tạp mà mọi hệ thống
streaming phải giải quyết. Trình phát video phải xử lý được các định dạng streaming hiện
đại (HLS, DASH) và tự động chọn chất lượng phù hợp với băng thông thực tế. Trang chủ cần
tổ chức nhiều khu vực nội dung khác nhau (banner, phim mới, series nổi bật, sự kiện thể
thao) với cơ chế tải bất đồng bộ đủ linh hoạt để không gây giật khi cuộn hoặc làm chậm
phản hồi giao diện. Hệ thống phải phân biệt rõ ràng giữa người xem thông thường và quản
trị viên thông qua cơ chế xác thực không trạng thái (stateless authentication) giữa ứng
dụng di động và máy chủ web. Mọi thao tác đều đi qua mạng di động — môi trường vốn không
ổn định về băng thông và độ trễ — đòi hỏi cơ chế caching, xử lý lỗi và tự phục hồi hiệu
quả. Cuối cùng, ứng dụng phải chạy ổn định trên nhiều dòng thiết bị Android khác nhau về
kích thước màn hình và cấu hình phần cứng, đồng thời tuân thủ vòng đời (lifecycle) chặt
chẽ của hệ điều hành.

Từ nhu cầu thực tế nêu trên, đồ án \textbf{CineFlow} được thực hiện với mục tiêu nghiên
cứu và xây dựng một ứng dụng xem phim trực tuyến hoàn chỉnh trên nền tảng Android, áp dụng
kiến trúc client--server và các kỹ thuật lập trình di động chuẩn mực để giải quyết các
thách thức kỹ thuật đó trong điều kiện thực tế.


% ────────────────────────────────────────────────────────────
\section{Mục tiêu và phạm vi đồ án}
\label{sec:muc_tieu_pham_vi}

\subsection{Mục tiêu}
\label{subsec:muc_tieu}

Mục tiêu của đồ án là xây dựng một ứng dụng xem phim trực tuyến hoàn chỉnh trên nền tảng
Android — từ phía client đến phía server — với các chức năng cốt lõi gồm: quản lý tài
khoản, khám phá nội dung, xem phim lẻ và series, xem video ngắn (shorts), theo dõi sự
kiện thể thao và quản trị nội dung. Bên cạnh việc hoàn thiện sản phẩm về mặt chức năng,
nhóm tập trung giải quyết các bài toán kỹ thuật thực tế trong lập trình di động:

\begin{itemize}
    \item Áp dụng kiến trúc client--server với REST API — ứng dụng Android giao tiếp với
    máy chủ Spring Boot qua giao thức HTTP, dữ liệu trao đổi dưới dạng JSON.

    \item Tổ chức ứng dụng Android theo mô hình phân lớp (Layered Architecture) kết hợp
    Repository Pattern, tách biệt rõ ràng tầng giao diện, tầng dữ liệu và tầng giao tiếp
    mạng.

    \item Xây dựng luồng xác thực không trạng thái dựa trên JSON Web Token (JWT) kết hợp
    refresh token rotation, cho phép duy trì phiên đăng nhập bền vững giữa các lần mở ứng
    dụng và tự động gia hạn token khi hết hạn.

    \item Tích hợp trình phát video chuyên dụng (Media3 ExoPlayer) hỗ trợ định dạng
    streaming hiện đại, có khả năng phát phim lẻ, series nhiều tập và sự kiện trực tiếp
    trong cùng một ứng dụng.

    \item Thiết kế hệ thống phân quyền hai vai trò (User và Admin) trong cùng một ứng
    dụng, kèm Admin Panel cho phép quản lý nội dung (CRUD phim, tập phim, danh mục, người
    dùng) và xem thống kê.
\end{itemize}

\subsection{Phạm vi}
\label{subsec:pham_vi}

Đồ án bao gồm hai nhóm người dùng chính:

\begin{itemize}
    \item \textbf{Người xem (User):} Đăng ký/đăng nhập, duyệt trang chủ, xem phim lẻ,
    phim series, video ngắn, nội dung thể thao; quản lý danh sách yêu thích, lịch sử xem,
    thông tin tài khoản và cài đặt ứng dụng.

    \item \textbf{Quản trị viên (Admin):} Ngoài các chức năng của User, có thêm khả năng
    quản lý nội dung hệ thống qua Admin Panel (tạo/sửa/xoá phim, tập phim, danh mục,
    người dùng) và xem thống kê hệ thống.
\end{itemize}

Về mặt nền tảng, đồ án triển khai trên \textbf{Android} (minSdk 24, targetSdk 36).
Ứng dụng cho iOS, phiên bản Web và Smart TV nằm ngoài phạm vi của đồ án này.


% ────────────────────────────────────────────────────────────
\section{Tổng quan kiến trúc hệ thống}
\label{sec:tong_quan_kien_truc}

CineFlow được thiết kế theo kiến trúc client--server, trong đó ứng dụng Android đóng vai
trò phía client và máy chủ Spring Boot đóng vai trò phía server. Hai phía giao tiếp qua
REST API trên nền giao thức HTTP, dữ liệu trao đổi dưới định dạng JSON.

% ── 3.1 Phía client (Android) ──────────────────────────────
\subsection{Kiến trúc phía client}
\label{subsec:kien_truc_client}

Ứng dụng Android được tổ chức theo mô hình phân lớp, gồm các tầng chính:

\begin{itemize}
    \item \textbf{Tầng giao diện (UI):} Chứa các Activity, Fragment và Adapter chịu trách
    nhiệm hiển thị dữ liệu và xử lý tương tác người dùng. Mọi Fragment đều kế thừa từ một
    lớp cơ sở chung với luồng khởi tạo chuẩn hoá (khai báo layout → khởi tạo view → tải
    dữ liệu).

    \item \textbf{Tầng Repository:} Đóng vai trò trung gian duy nhất giữa tầng giao diện
    và các nguồn dữ liệu. Mỗi Repository cung cấp dữ liệu dưới dạng LiveData, cho phép
    giao diện quan sát và cập nhật tự động, đồng bộ với vòng đời của Activity/Fragment.

    \item \textbf{Tầng mạng (Network):} Chịu trách nhiệm giao tiếp HTTP với máy chủ,
    được xây dựng trên nền Volley kết hợp OkHttp làm HTTP backend. Tầng này tích hợp bộ
    chặn xác thực tự động đính kèm JWT vào mỗi request, và cơ chế tự động gia hạn token
    khi nhận phản hồi 401 từ máy chủ.

    \item \textbf{Tầng xác thực (Auth):} Quản lý phiên đăng nhập, lưu trữ access token,
    refresh token và thông tin người dùng vào bộ nhớ riêng của ứng dụng. Cung cấp phương
    thức kiểm tra trạng thái đăng nhập và vai trò người dùng.

    \item \textbf{Tầng lưu trữ nội bộ (Local):} Sử dụng Room Database để cache dữ liệu
    trang chủ, giảm số lượt gọi API khi người dùng mở lại ứng dụng.
\end{itemize}

% ── 3.2 Phía server (Spring Boot) ──────────────────────────
\subsection{Kiến trúc phía server}
\label{subsec:kien_truc_server}

Phía backend được xây dựng trên Spring Boot 3.5 (Java 17), tổ chức theo kiến trúc phân
lớp ba tầng cổ điển:

\begin{itemize}
    \item \textbf{Tầng Controller:} Mở các endpoint REST với tiền tố \texttt{/api/v1},
    nhận và validate dữ liệu vào.

    \item \textbf{Tầng Service:} Chứa logic nghiệp vụ, quản lý transaction và tích hợp
    với các dịch vụ bên ngoài (gửi email, gọi API bên ngoài, thao tác lưu trữ đối tượng).

    \item \textbf{Tầng Repository:} Truy cập cơ sở dữ liệu PostgreSQL qua Spring Data JPA,
    cung cấp các thao tác CRUD và truy vấn tuỳ chỉnh.
\end{itemize}


% ── 3.3 Luồng giao tiếp chính ───────────────────────────────
\subsection{Luồng giao tiếp chính}
\label{subsec:luong_giao_tiep}

Hình~\ref{fig:arch_overview} mô tả tổng quan kiến trúc và luồng giao tiếp chính giữa
các thành phần của CineFlow.

\begin{figure}[H]
    \centering
    \includegraphics[width=0.85\textwidth]{outputs/system_arch}
    \caption{Tổng quan kiến trúc hệ thống CineFlow}
    \label{fig:arch_overview}
\end{figure}

Các luồng giao tiếp chính bao gồm:

\begin{itemize}
    \item \textbf{Luồng xác thực:} Người dùng đăng nhập → máy chủ trả về access token
    (2 phút) + refresh token (7 ngày, lưu DB). Client lưu token vào bộ nhớ riêng, tự động
    đính kèm access token vào mỗi request. Khi access token hết hạn (nhận 401), cơ chế
    xác thực token tự động gọi endpoint refresh-token để lấy cặp token mới, sau đó thử
    lại request ban đầu.

    \item \textbf{Luồng tải nội dung:} Fragment gọi Repository → Repository gọi dịch vụ
    API qua Volley → request được xếp hàng và gửi đi → dữ liệu JSON trả về được parse
    bằng Gson → kết quả phát qua LiveData tới giao diện.

    \item \textbf{Luồng phát video:} Người dùng chọn tập phim → Activity phát video khởi
    tạo ExoPlayer với URL video → ExoPlayer tự động chọn chất lượng phù hợp băng thông
    (adaptive bitrate) và phát nội dung.

    \item \textbf{Luồng quản trị:} Admin truy cập Admin Panel → các màn hình quản trị gọi
    endpoint \texttt{/admin/**} → Spring Security kiểm tra vai trò ADMIN → Controller xử
    lý và trả về dữ liệu phân trang.
\end{itemize}



